\documentclass[12pt]{article}
\usepackage{amsmath}
\usepackage{amssymb}
\usepackage{geometry}
\usepackage{listings}
\usepackage{xcolor}
\usepackage{hyperref}

\geometry{margin=1in}

\title{LSTM Internal Mechanisms: Hidden State, Cell State, and Gates}
\author{Detailed Explanation for DDoS Detection Model}
\date{}

\begin{document}

\maketitle

\section{Overview}

The LSTM model in the code uses two LSTM layers:
\begin{lstlisting}[language=Python, basicstyle=\ttfamily]
LSTM(64, return_sequences=True, input_shape=(num_features, 1))
LSTM(32, return_sequences=False)
Dense(num_classes, activation='softmax')
\end{lstlisting}

\section{Core Components}

\subsection{Hidden State ($h_t$)}

\textbf{Definition:} The hidden state is a vector that carries information from previous timesteps and represents the "short-term memory" of the LSTM.

\textbf{Mathematical Representation:}
\[
h_t \in \mathbb{R}^d
\]
where $d$ is the number of LSTM units (64 for Layer 1, 32 for Layer 2).

\textbf{In Your Code:}
\begin{itemize}
    \item \textbf{Layer 1:} $h_t^{(1)} \in \mathbb{R}^{64}$ for each timestep $t$
    \item \textbf{Layer 2:} $h_t^{(2)} \in \mathbb{R}^{32}$ for each timestep $t$
\end{itemize}

\textbf{Flow Through Timesteps:}
\[
h_0 \xrightarrow{x_1} h_1 \xrightarrow{x_2} h_2 \xrightarrow{x_3} h_3 \xrightarrow{\cdots} h_T
\]

where $T = \text{num\_features}$ (typically 20 in your dataset).

\textbf{Output Behavior:}
\begin{itemize}
    \item \textbf{Layer 1:} $\text{return\_sequences=True}$ $\Rightarrow$ outputs all $h_1^{(1)}, h_2^{(1)}, \ldots, h_T^{(1)}$
    \item \textbf{Layer 2:} $\text{return\_sequences=False}$ $\Rightarrow$ outputs only final $h_T^{(2)}$
\end{itemize}

\subsection{Cell State ($c_t$)}

\textbf{Definition:} The cell state is the "long-term memory" that flows through the LSTM, storing information across many timesteps.

\textbf{Mathematical Representation:}
\[
c_t \in \mathbb{R}^d
\]

\textbf{Initialization:}
\[
c_0 = \mathbf{0} \quad \text{(zero vector)}
\]

\textbf{Update Rule:}
\[
c_t = f_t \odot c_{t-1} + i_t \odot \tilde{c}_t
\]

where:
\begin{itemize}
    \item $f_t$ = forget gate (what to forget)
    \item $i_t$ = input gate (what to add)
    \item $\tilde{c}_t$ = candidate values (new information)
    \item $\odot$ = element-wise multiplication (Hadamard product)
\end{itemize}

\textbf{Key Properties:}
\begin{itemize}
    \item Flows through the entire sequence: $c_0 \rightarrow c_1 \rightarrow c_2 \rightarrow \cdots \rightarrow c_T$
    \item Modified by gates at each timestep
    \item Same dimensionality as hidden state
\end{itemize}

\section{LSTM Gates}

All gates use sigmoid activation: $\sigma(x) = \frac{1}{1 + e^{-x}} \in [0, 1]$

\subsection{Forget Gate ($f_t$)}

\textbf{Purpose:} Decide what information to forget from the cell state.

\textbf{Mathematical Formula:}
\[
f_t = \sigma(W_f \cdot [h_{t-1}, x_t] + b_f)
\]

where:
\begin{itemize}
    \item $W_f \in \mathbb{R}^{d \times (d + d_{in})}$ = forget gate weight matrix
    \item $b_f \in \mathbb{R}^d$ = forget gate bias vector
    \item $[h_{t-1}, x_t]$ = concatenation of previous hidden state and current input
    \item $d_{in}$ = input dimension (1 in your code: $\text{num\_features} \times 1$)
\end{itemize}

\textbf{Output Interpretation:}
\[
f_t[i] = \begin{cases}
\text{close to 1} & \text{keep information in } c_{t-1}[i] \\
\text{close to 0} & \text{forget information in } c_{t-1}[i]
\end{cases}
\]

\textbf{In Your Code:}
\begin{itemize}
    \item \textbf{Layer 1:} $f_t^{(1)} \in \mathbb{R}^{64}$, processes each feature sequentially
    \item \textbf{Layer 2:} $f_t^{(2)} \in \mathbb{R}^{32}$, processes hidden states from Layer 1
\end{itemize}

\textbf{Example:}
\[
f_t = [0.9, 0.1, 0.8, 0.95, \ldots]
\]
\begin{itemize}
    \item Dimension 1: Keep 90\% of old cell state
    \item Dimension 2: Keep 10\% of old cell state (forget 90\%)
    \item Dimension 3: Keep 80\% of old cell state
    \item Dimension 4: Keep 95\% of old cell state
\end{itemize}

\subsection{Input Gate ($i_t$) and Candidate Values ($\tilde{c}_t$)}

\textbf{Purpose:} Decide what new information to store in the cell state.

\textbf{Input Gate:}
\[
i_t = \sigma(W_i \cdot [h_{t-1}, x_t] + b_i)
\]

\textbf{Candidate Values:}
\[
\tilde{c}_t = \tanh(W_C \cdot [h_{t-1}, x_t] + b_C)
\]

where:
\begin{itemize}
    \item $W_i \in \mathbb{R}^{d \times (d + d_{in})}$ = input gate weight matrix
    \item $W_C \in \mathbb{R}^{d \times (d + d_{in})}$ = candidate weight matrix
    \item $b_i, b_C \in \mathbb{R}^d$ = bias vectors
    \item $\tanh(x) = \frac{e^x - e^{-x}}{e^x + e^{-x}} \in [-1, 1]$ = hyperbolic tangent
\end{itemize}

\textbf{Interpretation:}
\begin{itemize}
    \item $i_t[i]$ = how much of candidate value $\tilde{c}_t[i]$ to add
    \item $\tilde{c}_t[i]$ = new information value (between -1 and 1)
\end{itemize}

\textbf{Cell State Update:}
\[
c_t = f_t \odot c_{t-1} + i_t \odot \tilde{c}_t
\]

This combines:
\begin{enumerate}
    \item \textbf{Forgetting:} $f_t \odot c_{t-1}$ (what to keep from old state)
    \item \textbf{Adding:} $i_t \odot \tilde{c}_t$ (what new information to add)
\end{enumerate}

\textbf{Example:}
\[
\begin{aligned}
c_{t-1} &= [0.5, -0.3, 0.8, \ldots] \\
f_t &= [0.9, 0.2, 0.7, \ldots] \\
i_t &= [0.6, 0.8, 0.3, \ldots] \\
\tilde{c}_t &= [0.4, -0.5, 0.2, \ldots] \\
c_t &= [0.9, 0.2, 0.7, \ldots] \odot [0.5, -0.3, 0.8, \ldots] + [0.6, 0.8, 0.3, \ldots] \odot [0.4, -0.5, 0.2, \ldots] \\
&= [0.45, -0.06, 0.56, \ldots] + [0.24, -0.40, 0.06, \ldots] \\
&= [0.69, -0.46, 0.62, \ldots]
\end{aligned}
\]

\subsection{Output Gate ($o_t$)}

\textbf{Purpose:} Decide what parts of the cell state to output as the hidden state.

\textbf{Mathematical Formula:}
\[
o_t = \sigma(W_o \cdot [h_{t-1}, x_t] + b_o)
\]

where:
\begin{itemize}
    \item $W_o \in \mathbb{R}^{d \times (d + d_{in})}$ = output gate weight matrix
    \item $b_o \in \mathbb{R}^d$ = output gate bias vector
\end{itemize}

\textbf{Hidden State Computation:}
\[
h_t = o_t \odot \tanh(c_t)
\]

\textbf{Interpretation:}
\begin{itemize}
    \item $\tanh(c_t)$ = filtered cell state (scaled to $[-1, 1]$)
    \item $o_t$ = filter to decide what to output
    \item $h_t$ = final hidden state (what the LSTM "remembers" for next timestep)
\end{itemize}

\textbf{Example:}
\[
\begin{aligned}
c_t &= [0.69, -0.46, 0.62, \ldots] \\
\tanh(c_t) &= [0.60, -0.43, 0.55, \ldots] \\
o_t &= [0.8, 0.3, 0.9, \ldots] \\
h_t &= [0.8, 0.3, 0.9, \ldots] \odot [0.60, -0.43, 0.55, \ldots] \\
&= [0.48, -0.13, 0.50, \ldots]
\end{aligned}
\]

\section{Complete LSTM Computation at Timestep $t$}

\subsection{Step-by-Step Algorithm}

For each timestep $t$ (each feature in your code):

\textbf{Input:}
\begin{itemize}
    \item $x_t$ = current feature value (scalar in your case: $x_t \in \mathbb{R}$)
    \item $h_{t-1}$ = previous hidden state ($h_{t-1} \in \mathbb{R}^d$)
    \item $c_{t-1}$ = previous cell state ($c_{t-1} \in \mathbb{R}^d$)
\end{itemize}

\textbf{Step 1: Concatenate Input}
\[
z_t = [h_{t-1}, x_t] \in \mathbb{R}^{d + d_{in}}
\]

\textbf{Step 2: Compute Forget Gate}
\[
f_t = \sigma(W_f \cdot z_t + b_f)
\]

\textbf{Step 3: Compute Input Gate}
\[
i_t = \sigma(W_i \cdot z_t + b_i)
\]

\textbf{Step 4: Compute Candidate Values}
\[
\tilde{c}_t = \tanh(W_C \cdot z_t + b_C)
\]

\textbf{Step 5: Update Cell State}
\[
c_t = f_t \odot c_{t-1} + i_t \odot \tilde{c}_t
\]

\textbf{Step 6: Compute Output Gate}
\[
o_t = \sigma(W_o \cdot z_t + b_o)
\]

\textbf{Step 7: Compute Hidden State}
\[
h_t = o_t \odot \tanh(c_t)
\]

\textbf{Output:}
\begin{itemize}
    \item $h_t$ = new hidden state
    \item $c_t$ = new cell state
\end{itemize}

\subsection{Matrix Formulation}

All gates can be computed simultaneously:

\[
\begin{bmatrix}
f_t \\
i_t \\
o_t \\
\tilde{c}_t
\end{bmatrix}
=
\begin{bmatrix}
\sigma(W_f \cdot z_t + b_f) \\
\sigma(W_i \cdot z_t + b_i) \\
\sigma(W_o \cdot z_t + b_o) \\
\tanh(W_C \cdot z_t + b_C)
\end{bmatrix}
\]

Then:
\[
\begin{aligned}
c_t &= f_t \odot c_{t-1} + i_t \odot \tilde{c}_t \\
h_t &= o_t \odot \tanh(c_t)
\end{aligned}
\]

\section{In Your Specific Code}

\subsection{Layer 1: LSTM(64)}

\textbf{Input Shape:} $(N, T, 1)$ where:
\begin{itemize}
    \item $N$ = batch size
    \item $T$ = num\_features (typically 20)
    \item $1$ = feature dimension
\end{itemize}

\textbf{Processing:}
\[
\begin{aligned}
\text{For } t = 1, 2, \ldots, T: \\
&\quad x_t \in \mathbb{R} \quad \text{(feature } t \text{ value)} \\
&\quad h_{t-1}^{(1)} \in \mathbb{R}^{64} \quad \text{(previous hidden state)} \\
&\quad c_{t-1}^{(1)} \in \mathbb{R}^{64} \quad \text{(previous cell state)} \\
&\quad z_t = [h_{t-1}^{(1)}, x_t] \in \mathbb{R}^{65} \\
&\quad \text{Compute gates: } f_t^{(1)}, i_t^{(1)}, o_t^{(1)}, \tilde{c}_t^{(1)} \in \mathbb{R}^{64} \\
&\quad c_t^{(1)} = f_t^{(1)} \odot c_{t-1}^{(1)} + i_t^{(1)} \odot \tilde{c}_t^{(1)} \\
&\quad h_t^{(1)} = o_t^{(1)} \odot \tanh(c_t^{(1)})
\end{aligned}
\]

\textbf{Output:} Since $\text{return\_sequences=True}$:
\[
H^{(1)} = [h_1^{(1)}, h_2^{(1)}, \ldots, h_T^{(1)}] \in \mathbb{R}^{T \times 64}
\]

\subsection{Layer 2: LSTM(32)}

\textbf{Input:} $H^{(1)} \in \mathbb{R}^{T \times 64}$ (sequence from Layer 1)

\textbf{Processing:}
\[
\begin{aligned}
\text{For } t = 1, 2, \ldots, T: \\
&\quad h_t^{(1)} \in \mathbb{R}^{64} \quad \text{(hidden state from Layer 1)} \\
&\quad h_{t-1}^{(2)} \in \mathbb{R}^{32} \quad \text{(previous hidden state)} \\
&\quad c_{t-1}^{(2)} \in \mathbb{R}^{32} \quad \text{(previous cell state)} \\
&\quad z_t = [h_{t-1}^{(2)}, h_t^{(1)}] \in \mathbb{R}^{96} \quad \text{(concatenate 32 + 64)} \\
&\quad \text{Compute gates: } f_t^{(2)}, i_t^{(2)}, o_t^{(2)}, \tilde{c}_t^{(2)} \in \mathbb{R}^{32} \\
&\quad c_t^{(2)} = f_t^{(2)} \odot c_{t-1}^{(2)} + i_t^{(2)} \odot \tilde{c}_t^{(2)} \\
&\quad h_t^{(2)} = o_t^{(2)} \odot \tanh(c_t^{(2)})
\end{aligned}
\]

\textbf{Output:} Since $\text{return\_sequences=False}$:
\[
h_T^{(2)} \in \mathbb{R}^{32} \quad \text{(only final hidden state)}
\]

\subsection{Final Classification}

\[
\begin{aligned}
h_T^{(2)} &\in \mathbb{R}^{32} \quad \text{(final hidden state)} \\
\hat{y} &= \text{softmax}(W_{out} \cdot h_T^{(2)} + b_{out}) \\
\hat{y} &\in \mathbb{R}^{C} \quad \text{(class probabilities)}
\end{aligned}
\]

where $C = \text{num\_classes}$ (2 for binary classification).

\section{Parameter Count}

\subsection{Layer 1: LSTM(64)}

For each gate, weight matrix size: $W \in \mathbb{R}^{64 \times 65}$ (64 units, 65 inputs: 64 from $h_{t-1}$ + 1 from $x_t$)

\[
\begin{aligned}
\text{Forget Gate:} &\quad W_f: 64 \times 65 = 4,160 \text{ params}, \quad b_f: 64 \text{ params} \\
\text{Input Gate:} &\quad W_i: 64 \times 65 = 4,160 \text{ params}, \quad b_i: 64 \text{ params} \\
\text{Output Gate:} &\quad W_o: 64 \times 65 = 4,160 \text{ params}, \quad b_o: 64 \text{ params} \\
\text{Candidate:} &\quad W_C: 64 \times 65 = 4,160 \text{ params}, \quad b_C: 64 \text{ params} \\
\text{Total Layer 1:} &\quad 4 \times (4,160 + 64) = 16,896 \text{ parameters}
\end{aligned}
\]

\subsection{Layer 2: LSTM(32)}

Weight matrix size: $W \in \mathbb{R}^{32 \times 96}$ (32 units, 96 inputs: 32 from $h_{t-1}^{(2)}$ + 64 from $h_t^{(1)}$)

\[
\begin{aligned}
\text{Forget Gate:} &\quad W_f: 32 \times 96 = 3,072 \text{ params}, \quad b_f: 32 \text{ params} \\
\text{Input Gate:} &\quad W_i: 32 \times 96 = 3,072 \text{ params}, \quad b_i: 32 \text{ params} \\
\text{Output Gate:} &\quad W_o: 32 \times 96 = 3,072 \text{ params}, \quad b_o: 32 \text{ params} \\
\text{Candidate:} &\quad W_C: 32 \times 96 = 3,072 \text{ params}, \quad b_C: 32 \text{ params} \\
\text{Total Layer 2:} &\quad 4 \times (3,072 + 32) = 12,416 \text{ parameters}
\end{aligned}
\]

\subsection{Output Layer}

\[
\text{Dense: } 32 \times C + C \text{ parameters}
\]

\subsection{Total Parameters}

\[
\text{Total} = 16,896 + 12,416 + (32C + C) = 29,312 + 33C
\]

For binary classification ($C = 2$): $\text{Total} = 29,378$ parameters

\section{Key Differences: Hidden State vs Cell State}

\begin{table}[h]
\centering
\begin{tabular}{|l|l|l|}
\hline
\textbf{Aspect} & \textbf{Hidden State ($h_t$)} & \textbf{Cell State ($c_t$)} \\
\hline
Purpose & Short-term memory, output & Long-term memory, storage \\
Visibility & Exposed to next layer & Internal to LSTM \\
Update & Filtered cell state & Direct gate operations \\
Range & $[-1, 1]$ (via $\tanh$) & Unbounded (before $\tanh$) \\
Output & Used for predictions & Used internally \\
\hline
\end{tabular}
\caption{Comparison of Hidden State and Cell State}
\end{table}

\section{Summary}

The LSTM processes your tabular features sequentially:
\begin{enumerate}
    \item Each feature becomes a timestep
    \item At each timestep, gates control information flow
    \item Cell state stores long-term information
    \item Hidden state represents what to output
    \item Final hidden state goes to classification layer
\end{enumerate}

All gate computations are handled automatically by TensorFlow's \texttt{LSTM()} layer—you don't need to implement them manually.

\end{document}

